% KI-Übertragung in heutiges Deutsch aus der originalsprachlichen LaTeX-Fassung: gedruckte S. 252--260 und folgendes unnummeriertes Titelblatt [261].
% Hauptzeuge, PDF-Seiten 259--268; Google-Digitalisat, PDF-Seiten 263--272.
\begingroup
\clearpage
\noindent252\hfill\textsc{Briefwechsel.}\hfill\mbox{}\par\medskip
\noindent wie das zu machen ist. Im gegenwärtigen Fall muss man darauf ausgehen, die Bedingungsgleichung
\[
\alpha\mathrm d10+\beta\mathrm d12+\gamma\mathrm d13+\delta\mathrm d20+\varepsilon\mathrm d21+\zeta\mathrm d23+\eta\mathrm d30+\theta\mathrm d31+\varkappa\mathrm d32=\lambda
\]
so abzuändern, dass alle Korrekturen darin sind; das ist nun sehr leicht getan, man braucht nur
\[
\begin{aligned}
x(\mathrm d01-\mathrm d10+\mathrm d12-\mathrm d21+\mathrm d20-\mathrm d02)&=0\\
y(\mathrm d01-\mathrm d10+\mathrm d13-\mathrm d31+\mathrm d30-\mathrm d03)&=0\\
z(\mathrm d12-\mathrm d21+\mathrm d23-\mathrm d32+\mathrm d31-\mathrm d13)&=0
\end{aligned}
\]
hinzuaddieren, indem man $x$, $y$, $z$ nach Gefallen wählt. Aber so sind wir noch um nichts besser dran, wenn wir nicht wissen, wie wir wählen sollen. Ich sage: Wählt $x$, $y$, $z$ so, dass die Summe der 12 Quadrate der Koeffizienten in der entstehenden Gleichung ein Minimum bildet, also
\[
(x+y)^2+x^2+y^2+(\alpha-x-y)^2+(\beta+x+z)^2+\text{etc.}=\text{Minimum},
\]
woraus Sie $x$, $y$, $z$ mit leichter Mühe bestimmen.

Es lässt sich beweisen (freilich wird etwas künstliche Rechnung für diesen Beweis erfordert):

I. Dass man, wenn man dieses Geschäft viermal ausführte, nämlich zweitens ausgehend von $\frac{10}{12}.\allowbreak \frac{12}{13}.\allowbreak \frac{13}{10}=1$, dann von $\frac{20}{21}.\allowbreak \frac{21}{23}.\allowbreak \frac{23}{20}=1$ und viertens von $\frac{30}{31}.\allowbreak \frac{31}{32}.\allowbreak \frac{32}{30}=1$, 4 Endgleichungen erhält, die, genau besehen, untereinander identisch sind, d.\,h. die sämtlichen 13 Teile (den absoluten mitgezählt) sind in allen proportional. Der Rat, die erste, auf das größte Dreieck gegründete Entwicklung [zu nehmen], hat bloß den Zweck, alles in den größten Zahlen zu erhalten, die wirklich jeweils den Flächen der $\Delta\Delta$ $123$, $023$, $013$, $012$ proportional sind. Ich pflege übrigens zur Sicherheit alle 4 zu entwickeln.

Bei der Form \tikz[baseline=-.4ex,x=.027cm,y=.027cm]{\coordinate (a) at (0,0);\coordinate (b) at (1,10);\coordinate (c) at (13,11);\coordinate (d) at (16,0);\draw[line width=.8pt] (a)--(b)--(c)--(d)--cycle;\draw[line width=.3pt] (a)--(c) (b)--(d);} geben zwei zueinander addiert dieselbe Summe wie die beiden anderen; bei der Form \tikz[baseline=-.4ex,x=.027cm,y=.027cm]{\coordinate (a) at (0,0);\coordinate (b) at (8,11);\coordinate (c) at (20,0);\coordinate (o) at (8,5);\coordinate (t) at (8,8.2);\draw[line width=.8pt] (a)--(b)--(c)--cycle;\draw[line width=.3pt] (a)--(o)--(t) (o)--(c);} ist eine die Summe der drei anderen. Ich addiere alle 4, natürlich so gefasst, dass eine \textls[100]{absolute} Addition stattfindet.

So wird deutlich, dass der Symmetrie ihr volles Recht widerfahren ist. Übrigens lässt sich die Arbeit hierbei noch abkürzen. Darauf gehe ich nicht ein, da sie im Vergleich zum ganzen Geschäft jedenfalls ganz unbedeutend ist.

\clearpage
\noindent\mbox{}\hfill\textsc{Ausgleichung einfacher Figuren.}\hfill253\par\medskip
II. Aber es geht nicht bloß um die Symmetrie. Existierte nur ein solches Viereck, so würde man hier den großen Vorteil gewinnen, dass die Ausgleichung in dieser Form die Ausgleichung der 4 (von denen nur 3 voneinander unabhängig sind) Bedingungsgleichungen [der 2. Art], die dieses Viereck darbietet, gar nicht stört. Gibt es mehrere solche Vierecke, die voneinander getrennt sind (keine Seite gemeinsam haben), wie es in meinem System mehrfach der Fall war, so bleibt derselbe Vorteil für alle diese bestehen. Es werden dann immer nur die Winkelsummen in den \textls[100]{angrenzenden} $\Delta\Delta$ gestört. Aber ich bin überzeugt, dass man selbst bei einem so verschränkten System wie dem Ihren auf diesem Weg eine \textls[100]{sehr bedeutend schnellere} Ausgleichung erreichen würde.

Ich glaube, es wird Ihnen Vergnügen machen, dieses Verfahren einmal auch nur auf eines Ihrer Vierecke anzuwenden. Bei 5-Ecken, 6-Ecken usw. ist es übrigens ganz analog.

Ich hatte dieses Verfahren wirklich beinahe vergessen, sodass ich es erst beim Durchsehen alter Rechnungen von 1826 wiederfand. Denn bei den späteren $\Delta\Delta$, wie auch dieses Jahr bei den Messungen im Bremischen, gehe ich etwas anders vor, da sie den großen Zeitaufwand nicht verdienen. Ich gleiche zunächst nur die Winkelsummen genau aus, worin ich eine solche Fertigkeit habe, dass ich z.\,B. für die sehr verschränkte Bremer Messung nur eine Stunde brauche. Dann gehe ich zu den Bedingungsgleichungen der 3. Art (Seitenverhältnisse) über. Ich behandle sie so, dass sie genau erfüllt werden und jene Winkelsummen dabei überhaupt \textls[100]{nicht} wieder gestört werden. Dies ergibt zwar nicht das \textls[100]{absolute Minimum} der Quadratsumme, bleibt aber jedenfalls nicht weit dahinter zurück und bringt alles in vollkommene Übereinstimmung. Bei den westfälischen Messungen habe ich das Verfahren wohl so abgeändert, dass ich das Geschäft mehrere Male durchführte. Das lässt sich so einrichten, dass man der absoluten kleinsten Quadratsumme immer näher kommt, je öfter man es wiederholt. Es ist jedoch ganz unmöglich, dieses Verfahren in aller Kürze zu beschreiben.

\clearpage
\noindent254\hfill\textsc{Nachlass.}\hfill\mbox{}\par\medskip
\begin{center}[4.]\\[5pt]Zur Ausgleichung der Winkel im Viereck.\end{center}
Es seien die 4 Punkte $a$, $b$, $c$, $d$;
\[
\begin{array}{rlllll}
A&\text{doppelter}&\text{Inhalt}&\text{des}&\text{Dreiecks}&bcd\\
{[}B&\text{»}&\text{»}&\text{»}&\text{»}&cda\\
C&\text{»}&\text{»}&\text{»}&\text{»}&dab\\
D&\text{»}&\text{»}&\text{»}&\text{»}&abc].
\end{array}
\]
Der Winkel $bcd$ werde mit $Ac$ bezeichnet, und so die übrigen.

Man setze
\[
\begin{aligned}
\log\tfrac{\sin Bc.\sin Cd.\sin Db}{\sin Bd.\sin Cb.\sin Dc}&=\alpha\\[4pt]
\log\tfrac{\sin Cd.\sin Da.\sin Ac}{\sin Ca.\sin Dc.\sin Ad}&=\beta\\[4pt]
\log\tfrac{\sin Da.\sin Ab.\sin Bd}{\sin Db.\sin Ad.\sin Ba}&=\gamma\\[4pt]
\llap{\raisebox{-.5ex}{\text{\tiny.}}\hspace{3.5cm}}\log\tfrac{\sin Ab.\sin Bc.\sin Ca}{\sin Ac.\sin Ba.\sin Cb}&=\delta.
\end{aligned}
\]
Man hat dann [wenn die Winkelsummen in den Dreiecken bereits ausgeglichen sind:]
\[
\frac\alpha A=\frac\beta B=\frac\gamma C=\frac\delta D=\mu,
\]
\noindent[und näherungsweise nach Art. 1, indem man die Richtungsverbesserungen $C^{01}$ und $C^{10}$ usw., die jetzt mit $\partial ab$ und $\partial ba$ usw. bezeichnet werden, gleichsetzt und gleichzeitig $A$ für $-T^0$, $B$ für $T'$, $C$ für $-T''$ und $D$ für $T'''$ schreibt,]
\[
M\mu=\frac{ab^2}{CD}.\partial ab-\frac{ac^2}{BD}.\partial ac+\frac{ad^2}{BC}.\partial ad+\frac{bc^2}{AD}.\partial bc-\frac{bd^2}{AC}.\partial bd+\frac{cd^2}{AB}.\partial cd;
\]
\[
M=\frac{206265}{k},\qquad k=\text{Modul der Briggs'schen Logarithmen}
\]
\noindent\hspace*{.7cm}$\log M=5{,}676\,6408$.
\begin{center}Beispiel.\end{center}
\begin{center}
\begin{tabular}{rcccc}
Viereck:&Deister,&Lichtenberg,&Garssen,&Falkenberg.\\
&$d$&$c$&$b$&$a$
\end{tabular}
\end{center}
[Die Ausgleichung der Winkelsummen der 4 Dreiecke des Vierecks in der Ebene hatte die folgenden Werte ergeben.]

\clearpage
\noindent\mbox{}\hfill\textsc{Ausgleichung einfacher Figuren.}\hfill255\par\medskip
\[
\setlength{\arraycolsep}{2pt}
\begin{array}{l|rrr|l|rrr|l|rrr|l|rrr}
[Ab&47^\circ&18'&46''{,}567&[Ba&34^\circ&33'&59''{,}066&[Ca&57^\circ&33'&17''{,}328&[Da&22^\circ&59'&18''{,}262\\
Ac&66&1&16{,}315&Bc&55&34&13{,}347&Cb&99&14&52{,}202&Db&146&33&38{,}770\\
Ad]&66&39&57{,}118&Bd]&89&51&47{,}587&Cd]&23&11&50{,}470&Dc]&10&27&2{,}968\\ \hline
&180&0&0{,}000&&180&0&0{,}000&&180&0&0{,}000&&180&0&0{,}000
\end{array}
\]
[Die Näherungswerte der Logarithmen der Seiten sind:
\[
\begin{array}{l@{\;}r@{\qquad}l@{\;}r}
ab\ .\ .\ .&4{,}44961&bc\ .\ .\ .&4{,}78266\\
ac\ .\ .\ .&4{,}93218&bd\ .\ .\ .&4{,}78052\\
ad\ .\ .\ .&4{,}84854&cd\ .\ .\ .&4{,}68605.
\end{array}
\]
Damit ergibt sich:]
\[
\begin{array}{l@{\;}r@{\qquad}l@{\;}r@{\qquad}l@{\;}r}
\log A=&9{,}42951&\sin Cd\ .\ .\ .&9{,}595\,3854&\sin Ca\ .\ .\ .&9{,}926\,2936\\
\log B=&9{,}53459&\sin Da\ .\ .\ .&9{,}591\,6709&\sin Dc\ .\ .\ .&9{,}258\,6170\\
\log C=&9{,}22445&\sin Ac\ .\ .\ .&9{,}960\,8017&\sin Ad\ .\ .\ .&9{,}962\,9423\\ \cline{4-4}\cline{6-6}
\log D=&8{,}97346&&9{,}147\,8580&&9{,}147\,8529
\end{array}
\]
\noindent\makebox[\linewidth][c]{\hspace{.27\linewidth}$[\beta=0{,}000\,0051]$}\par\medskip
\[
\setlength{\arraycolsep}{3pt}
\begin{array}{l@{}r|@{\quad}l@{}r@{\quad}l@{}r@{\quad}l@{}r}
[\beta.\ .\ .\ .\ .\ &4{,}70757&[ab^2\ .\ .\ .\ &8{,}89922&[ac^2\ .\ .\ .\ .\ &\rlap{\raisebox{1.5ex}{\text{\tiny.}}}9{,}86436&[ad^2\ .\ .\ .\ .\ &9{,}69708\\
B\ .\ .\ .\ .\ &\underline{9{,}53459}&CD\ .\ .\ .\ &8{,}19791&BD\ .\ .\ .\ &8{,}50805&BC\ .\ .\ .\ .\ &8{,}75904\\ \cline{4-4}\cline{6-6}\cline{8-8}
\mu\ .\ .\ .\ .\ &5{,}17298&&0{,}70131&&1{,}35631&&0{,}93804\\[5pt]
M\ .\ .\ .\ .\ &\underline{5{,}67664}&bc^2.\ .\ .\ .\ &9{,}56532&bd^2\ .\ .\ .\ .\ &9{,}56104&cd^2\ .\ .\ .\ .\ &9{,}37210\\ 
M\mu.\ .\ ]&0{,}84962&AD\ .\ .\ ]&8{,}40297&AC\ .\ .\ ]&8{,}65396&AB\ .\ .\ .\ ]&8{,}96410\\ \cline{4-4}\cline{6-6}\cline{8-8}
&&&1{,}16235&&0{,}90708&&0{,}40800
\end{array}
\]
\[
\begin{aligned}
7''{,}073={}&5{,}027\partial ab-22{,}715\partial ac+8{,}670\partial ad+14{,}533\partial bc\\
&\hspace{5.5cm}-8{,}074\partial bd+2{,}559\partial cd.
\end{aligned}
\]
[Hieraus erhält man für die Verbesserungen
\[
\partial ab=5{,}027\varkappa,\qquad\partial ac=-22{,}715\varkappa,\quad\text{usw.,}
\]
\noindent wo
\[
899{,}35\varkappa=7{,}073,
\]
\noindent also
\[
\log\varkappa=7{,}89567-10
\]
\noindent ist.]

\clearpage
\noindent256\hfill\textsc{Nachlass.}\hfill\mbox{}\par\medskip
\[
\begin{array}{r@{\;}c@{\;}r}
\partial ab&=&+0''{,}0395\\
\partial ac&=&-0{,}1786\\
\partial ad&=&+0{,}0682.\\
\partial bc&=&+0{,}1143\\
\partial bd&=&-0{,}0635\\
\partial cd&=&+0{,}0201.
\end{array}
\]
[Die Winkelverbesserungen sind daher
\[
\begin{aligned}
\partial Ab&=-\partial bc+\partial bd=-0''{,}178\\
\partial Ac&=-\partial cd+\partial bc=+0{,}094\\
\partial Ad&=-\partial bd+\partial cd=+0{,}084
\end{aligned}
\]
\begin{center}usw.\end{center}
Mit dem gegebenen Wert der Seite $cd$ (Lichtenberg---Deister) als Ausgangswert für die Seitenberechnung ergibt sich mithin die folgende Zusammenstellung:]
\begin{center}
\renewcommand{\arraystretch}{1.08}
\begin{tabular}{|l|r@{\quad}r@{\quad}r|r|r|}
\hline
\multicolumn{1}{|c|}{[ Station}&\multicolumn{3}{c|}{Winkel}&\multicolumn{1}{c|}{$\log\sin$}&\begin{tabular}{c}Log.\\der Seiten\end{tabular} ]\\ \hline\hline
Garssen&$47^\circ$&$18'$&$46''{,}389$&$9{,}866\,3270$&$4{,}686\,0451$\\
Lichtenberg&66&1&$16{,}409$&$9{,}960\,8018$&$4{,}780\,5199$\\
Deister&66&39&$57{,}202$&$9{,}962\,9424$&$4{,}782\,6605$\\[9pt]
Falkenberg&34&33&$59{,}313$&$9{,}753\,8603$&$4{,}686\,0451$\\
Lichtenberg&55&34&$13{,}148$&$9{,}916\,3595$&$4{,}848\,5443$\\
Deister&89&51&$47{,}539$&$9{,}999\,9988$&$4{,}932\,1836$\\[9pt]
Falkenberg&57&33&$17{,}357$&$9{,}926\,2937$&$4{,}780\,5199$\\
Garssen&99&14&$52{,}305$&$9{,}994\,3182$&$4{,}848\,5444$\\
Deister&23&11&$50{,}338$&$9{,}595\,3848$&$4{,}449\,6110$\\[9pt]
Falkenberg&22&59&$18{,}044$&$9{,}591\,6698$&$4{,}782\,6605$\\
Garssen&146&33&$38{,}695$&$9{,}741\,1930$&$4{,}932\,1837$\\
Lichtenberg&10&27&$3{,}261$&$9{,}258\,6204$&$4{,}449\,6111$\\[5pt]
\end{tabular}
\end{center}

\clearpage
\noindent\mbox{}\hfill\textsc{Ausgleichung einfacher Figuren.}\hfill257\rlap{\hspace{.4em}\raisebox{1.3ex}{\tiny.}}\par\medskip
\begin{center}[5.]\hspace{.5em}.\\[5pt]Viereck zwischen 4 Punkten 1. 2. 3. 4.\end{center}
Die Bedingungsgleichung wird so gebildet: Das Produkt der Lote auf die Seite $12$ aus den Punkten $3$ und $4$ sei $p'p'$ usw. Bezeichnet man nun die Richtung der Seite $12$ mit $a$ [ihre Verbesserung mit $\mathrm da=\mathrm d12$], setzt $p'p'=A$ und verfährt ebenso bei den übrigen, so gilt:
\[
0=\frac{\mathrm da}A+\frac{\mathrm db}B+\cdots
\]
\noindent[oder]
\[
\begin{aligned}
0={}&\mathrm d12-\tfrac{\sin142\sin321}{\sin314\sin132}\,\mathrm d13+\tfrac{\sin213\sin421}{\sin314\sin142}\,\mathrm d14\\
&+\tfrac{\sin213\sin421}{\sin324\sin132}\,\mathrm d23-\tfrac{\sin214\sin321}{\sin324\sin142}\,\mathrm d24
+\left\{\begin{array}{cl}
&\tfrac{\sin214\sin321}{\sin143\sin432}\,\mathrm d34\\[5pt]
\text{oder}&\tfrac{\sin213\sin421}{\sin431\sin243}\,\mathrm d34.
\end{array}\right.
\end{aligned}
\]
\bigskip
\begin{center}[6.]\\[5pt][Ausgleichung eines Polygons.]\end{center}
Schließt sich eine Anzahl von $n$ Dreiecken um einen einzigen Punkt zusammen und ist
\[
\begin{array}{rlllllll}
A&\text{Korrektur}&\text{der Summe}&\text{der Winkel}&a,a',a''&\text{im}&\text{ersten}&\text{Dreieck}\\
B&\text{»}&\text{»}&\text{»}&b,b',b''&\text{»}&\text{zweiten}&\text{»}\\
C&\text{»}&\text{»}&\text{»}&c,c',c''&\text{»}&\text{dritten}&\text{»}
\end{array}
\]
\begin{center}usw.,\end{center}
und $S$ die Korrektur der Summe der Winkel $a+b+c+\text{usw.}$ [die um den Punkt herum liegen], so korrigiert man
\[
\begin{array}{r@{\quad}l@{\;}l}
a&\text{um}&\tfrac{n-1}{3n}A-\tfrac1{3n}B-\tfrac1{3n}C-\cdots+\tfrac1nS\\[6pt]
\left.\begin{array}{r}a'\\a''\end{array}\right\}&\text{jeden um}&\tfrac{2n+1}{6n}A+\tfrac1{6n}B+\tfrac1{6n}C+\cdots-\tfrac1{2n}S
\end{array}
\]
\begin{center}usw.\end{center}
[Vorausgesetzt ist hierbei, dass man nur die Winkelgleichungen, nicht
\vfill\noindent{\small IX.}\hfill33

\clearpage
\noindent258\hfill\textsc{Nachlass.}\hfill\mbox{}\par\medskip
\noindent aber die Seitengleichung, in Rücksicht zieht, und dass alle Winkelbeobachtungen gleiches Gewicht, 1, haben.]

Diese Vorschrift kann auch folgendermaßen formuliert werden.

1) Man verbessere zuerst jeden Winkel im ersten Dreieck um $\frac13A$, im zweiten um $\frac13B$, usw.

2) Sodann verbessere man nochmals die um den einen Punkt herum liegenden Winkel, sodass sie schließen.

3) Dann wende man noch die Hälfte dieser zweiten Korrektur mit entgegengesetztem Vorzeichen auf jeden der übrigen Winkel an.

Der mittlere noch zu befürchtende Fehler [eines jeden $a$, $b$, $c$, $\ldots$ ist gleich]
\[
m\sqrt{\frac{2n-2}{3n}}
\]
\noindent[und eines jeden $a'$, $b'$, $a''$, $b''$, $\ldots$ gleich]
\[
m\sqrt{\frac{4n-1}{6n}}.
\]
[$m$ ist der mittlere Fehler der Gewichtseinheit.]
\bigskip
\begin{center}[7.]\\[5pt]Gewicht von Höhenbestimmungen.\end{center}
Mit einem Anfangspunkt $P$ ist eine Reihe anderer [Punkte] $p$, $p'$, $p''$, $p'''$, $\ldots$ $p^{(n)}$ usf. so verbunden, dass die Höhenunterschiede $Pp$, $Pp'$, $pp'$, $pp''$, $p'p''$, $p'p'''$, $p''p'''$, $p''p^{\mathrm{IV}}$ usf., also die Unterschiede jeder Höhe zu den Höhen der beiden vorhergehenden und der beiden folgenden Punkte, mit gleicher Zuverlässigkeit beobachtet worden sind. Dann ist das Gewicht der Bestimmung des Höhenunterschieds zwischen $P$ und $p^{(n)}$ nach der Methode der kleinsten Quadrate
\[
\setlength{\arraycolsep}{7pt}
\begin{array}{r@{\;}cccccccccc}
n=&0&1&2&3&4&5&6&7&8&9\\[4pt]
\text{Gewicht:}&\frac11&\frac32&\frac88&\frac{21}{24}&\frac{55}{75}&\frac{144}{224}&\frac{377}{663}&\frac{987}{1932}&\frac{2584}{5576}&\frac{6765}{15950}.
\end{array}
\]
Die Zähler sind die Koeffizienten der Reihe aus der Entwicklung des Bruchs
\[
\frac1{1-3x+xx}
\]

\clearpage
\noindent\mbox{}\hfill\textsc{Ausgleichung einfacher Figuren.}\hfill259\par\medskip
\noindent und folglich in der Formel
\[
\tfrac15(3r^{n+1}-2r^n-2r^{-n}+3r^{-n-1})
\]
\noindent begriffen, wo
\[
\begin{gathered}
2r=3+\sqrt5\quad\text{oder}\quad r=2{,}618\,0340\\
\log r=0{,}417\,9752
\end{gathered}
\]
\noindent ist. Ebenso entstehen die Nenner aus der Entwicklung des Bruchs
\[
\tfrac{(1-x)^2}{(1+x)(1-3x+xx)^2}
\]
\noindent und werden also dargestellt durch die Formel
\[
\tfrac1{25}\{(3n+7)r^{n+1}-(2n+2)r^n-(2n+2)r^{-n}+(3n+7)r^{-n-1}+(-1)^n8\}.
\]
Der Grenzausdruck für das Gewicht wird
\[
\frac{25}{5n+5+4\sqrt5}.
\]
\begin{center}\rule{3cm}{.4pt}\end{center}
\bigskip
\begin{center}\textsc{Bemerkungen.}\end{center}
\begingroup\footnotesize
Die Notiz [1] und der größte Teil von [2] sowie die Notiz [7] befinden sich in demselben Notizbuch. [6] gehört einem anderen Notizbuch an. Das Schlussergebnis von [2] steht auf der letzten Seite in \textsc{Gauss}’ Exemplar von \textsc{Carnot}s Geometrie der Stellung, übersetzt von \textsc{Schumacher}. Erster Theil, 1810.

Wenn in einem Viereck mithilfe von 3 Winkelgleichungen die Seitengleichung so umgeformt wird, wie es unter [3] im Brief an \textsc{Gerling} vom 19. Januar 1840 angegeben ist, lässt sich leicht zeigen (vgl. II, S. 253), dass sich die durch die Ausgleichung des Vierecks ergebenden Verbesserungen aus den Verbesserungen zusammensetzen, die die Ausgleichung allein der 3 Winkelgleichungen erfordert, und den Verbesserungen, die die Ausgleichung allein der umgeformten Seitengleichung erfordert. Denn bildet man aus den Winkelgleichungen und der umgeformten Seitengleichung die Ausdrücke für die Korrelaten und damit die Normalgleichungen, so hängen die den Winkelgleichungen entsprechenden Normalgleichungen mit der Normalgleichung, die aus der umgeformten Seitengleichung entsteht, \textls[100]{nicht} zusammen. Vorausgesetzt ist dabei, dass die Faktoren, mit denen die Winkelgleichungen vor der Addition zur ursprünglichen Seitengleichung multipliziert wurden, so bestimmt werden, dass die Summe der Quadrate der Koeffizienten in der neu entstandenen Seitengleichung ein Minimum wird. Es ist hierbei nicht nötig, dass die konstanten Glieder der Winkelgleichungen, die zur Umformung der Seitengleichung benutzt werden, vorher durch Ausgleichung auf null gebracht werden. Ist aber Letzteres der
\par\endgroup
\vfill\noindent\hfill33\textsuperscript{$\ast$}\hspace*{2cm}

\clearpage
\noindent260\hfill\textsc{Bemerkungen.\quad Ausgleichung einfacher Figuren.}\hfill\mbox{}\par\medskip
\begingroup\footnotesize
\noindent Fall, so muss die Konstante der Seitengleichung natürlich mit den Werten berechnet werden, die die Ausgleichung der Winkelgleichungen geliefert hat.

Die Aufzeichnung unter [4] ist einem zur hannoverschen Gradmessung gehörenden Rechenheft entnommen. Die Seitengleichungen in der angegebenen Form wurden dort bei einer Ausgleichung der aus den Jahren 1821 bis 1823 stammenden Dreiecke bis zur Seite Hamburg-Timpenberg aufgestellt. Nachdem zuerst die Winkelsummen der Dreiecke ausgeglichen worden sind, erfolgt die Ausgleichung der abgekürzten Bedingungsgleichungen, die sich aus den Seitenverhältnissen ergeben. Die Winkelsummen der Dreiecke werden durch diese Ausgleichung nicht geändert. Auf dieses Verfahren bezieht sich wohl der Schluss des unter [3] mitgeteilten Briefs an \textsc{Gerling} vom 19. Januar 1840. Für numerische Rechnungen erhält man diese Seitengleichung ebenso bequem, wenn man die Koeffizienten ihrer Verbesserungen wie üblich entweder mithilfe logarithmischer Differenzen oder mithilfe der Kotangenswerte der zugehörigen Winkel bildet und dann die Verbesserung der Richtung $hi$ gleich der Verbesserung der Richtung $ih$ setzt.

Die erste Formel des Art. [5] befindet sich im selben Rechenheft wie [4], während die zweite Formel auf dem Vorsatzblatt einer Logarithmentafel eingetragen ist. Bezeichnet man das Produkt der Lote von den Punkten $3$ und $4$ auf die Seite $12$ mit $P_{12}$, das Produkt der Lote aus $2$ und $4$ auf $13$ mit $P_{13}$ usw. und bezeichnen $\mathrm d12$, $\mathrm d13$ usw. die Richtungsänderungen von $12$, $13$ usw., so lautet nach Art. [4] die Seitengleichung, wenn das Viereck ausgeglichen ist, in Näherungsform:
\[
\frac{\mathrm d12}{P_{12}}-\frac{\mathrm d13}{P_{13}}+\frac{\mathrm d14}{P_{14}}+\frac{\mathrm d23}{P_{23}}-\frac{\mathrm d24}{P_{24}}+\frac{\mathrm d34}{P_{34}}=0.
\]
Hieraus ergibt sich sofort die zweite Formel des Art. [5]. Im Original hat das Glied mit $\mathrm d14$ der letzteren das Minusvorzeichen, während die übrigen Glieder sämtlich positiv sind. \textsc{Gauss} hat dieser Formel jedoch hinzugefügt: »Vorbehaltlich der Vorzeichen.«

Zu der Notiz [7] sei noch folgendes bemerkt. Ist
\[
\llap{\raisebox{1ex}{\text{\tiny.}}\hspace{4cm}}-N_{i-1}+3N_i-N_{i+1}=0,
\]
\noindent wobei $N_0=0$ und $N_1=1$ ist, so wird der Kehrwert des Gewichts des Höhenunterschieds durch die Formel
\[
\frac{(5n+17)N_{n+1}-8N_n+(-1)^n8}{25N_{n+1}}
\]
\noindent erhalten, aus der, wegen $\displaystyle\lim_{n=\infty}\frac{N_n}{N_{n+1}}=\frac{3-\sqrt5}{2}$, als Grenzausdruck
\[
\frac{5n+5+4\sqrt5}{25}
\]
\noindent folgt.\hfill\textsc{Krüger.}\par
\endgroup
\vspace{1.6cm}
\begin{center}\rule{4cm}{.4pt}\end{center}

\clearpage
% Unnummeriertes Titelblatt nach der gedruckten S. 260 ohne sichtbare Seitenzahl oder Bogensignatur.
\null\vspace*{\stretch{4}}
\begin{center}{\LARGE STATIONSAUSGLEICHUNGEN.}\end{center}
\vspace*{\stretch{6}}
\endgroup
% Ende des unnummerierten Titelblatts [261].
